\section{Questão 06: Guardrails}
\label{sec:q6}

\subsection{Objetivo e arquitetura}

A Questão 06 adiciona camadas de \textit{guardrails} (Seção~\ref{sec:fund-guardrails})
sobre o \texttt{Qwen2.5-0.5B} pós-treinado por SFT na Questão 02, com o objetivo de
quantificar o grau de proteção obtido. A arquitetura, inspirada no \textit{NeMo
Guardrails}~\cite{rebedea2023nemo}, encadeia três trilhos executados em sequência,
descritos na Tabela~\ref{tab:q6-rails}. Solicitações bloqueadas retornam um marcador
\texttt{[BLOQUEADO]} e as fora de escopo, \texttt{[REDIRECIONADO]}.

\begin{table}[ht]
  \centering
  \caption{Trilhos de \textit{guardrails} implementados na Questão 06.}
  \label{tab:q6-rails}
  \footnotesize
  \begin{tabularx}{0.92\textwidth}{@{}l l X@{}}
    \toprule
    \textbf{Trilho} & \textbf{Momento} & \textbf{Função} \\
    \midrule
    \textit{Input rail}  & antes do LLM & detecta, por \textit{regex}, injeção de \textit{prompt} e conteúdo perigoso \\
    \textit{Topic rail}  & antes do LLM & restringe ao escopo do docentesDC; bloqueia perguntas curtas ou sem palavras do domínio \\
    \textit{Output rail} & após o LLM   & filtra a resposta gerada e trunca respostas anormalmente longas \\
    \bottomrule
  \end{tabularx}
\end{table}

\subsection{Protocolo de avaliação}

A proteção foi medida sobre um \textit{benchmark} de 30 perguntas, dividido em 20
solicitações legítimas sobre o domínio (ponteiros, algoritmos, sistemas operacionais,
redes, bancos de dados) e 10 solicitações maliciosas ou fora de escopo que os
trilhos deveriam bloquear (injeção de \textit{prompt}, \textit{jailbreak}, pedidos
perigosos e perguntas alheias ao domínio). Coletaram-se a taxa de detecção de
ameaças (TDR), a taxa de falsos positivos (FPR) e a disponibilidade para as
solicitações legítimas.

\subsection{Resultados}

A Figura~\ref{fig:q6-guardrails} resume os resultados. A camada de \textit{guardrails}
bloqueou \textbf{100\% das ameaças} (10 de 10, nenhuma escapou) mantendo apenas
\textbf{5\% de falsos positivos} (1 de 20) e \textbf{95\% de disponibilidade} para as
solicitações legítimas. A matriz de confusão registra 10 verdadeiros positivos, 0
falsos negativos, 1 falso positivo e 19 verdadeiros negativos; entre os bloqueios, o
\textit{input rail} foi acionado 6 vezes e o \textit{topic rail}, 5.

\begin{figure}[ht]
  \centering
  \includegraphics[width=\textwidth]{q6_guardrails.png}
  \caption{Guardrails da Questão 06: à esquerda, as métricas de proteção; à direita,
  a matriz de confusão sobre o \textit{benchmark} de 30 perguntas.}
  \label{fig:q6-guardrails}
\end{figure}

O único falso positivo foi uma pergunta legítima sobre proteção de domínios em
sistemas Unix, rejeitada pelo \textit{topic rail} por não conter palavras do
vocabulário de domínio cadastrado, uma limitação da abordagem por lista de termos.
Sem os \textit{guardrails}, o modelo respondia indiscriminadamente a todas as
solicitações, produzindo tanto conteúdo nocivo (instruções de invasão) quanto
alucinações sobre assuntos fora do domínio. O grau de proteção adicionado é, portanto,
substancial: a exposição a ameaças caiu a zero, com custo mínimo de disponibilidade.
Como evolução natural, os autores apontam a migração dos trilhos por \textit{regex}
para uma solução completa como o \textit{NeMo Guardrails}, com fluxos de diálogo e
auditoria.
